\documentclass[11pt,a4paper]{article}
\usepackage[margin=0.85in]{geometry}
\usepackage{graphicx}
\usepackage{booktabs}
\usepackage{hyperref}
\hypersetup{
  colorlinks=true,
  linkcolor=black,
  citecolor=black,
  urlcolor=blue,
  pdfborder={0 0 0}
}
\usepackage{amsmath,amssymb}
\usepackage{needspace}
\usepackage{caption}
\usepackage{float}
\usepackage{placeins}
\usepackage{setspace}
\usepackage{microtype}
\usepackage{fancyhdr}
\usepackage{titlesec}
\usepackage{enumitem}
\setstretch{1.12}
\graphicspath{{outputs/}}
\setlist[itemize]{leftmargin=1.4em,itemsep=0.25em,topsep=0.35em}

\captionsetup{font=small,labelfont=bf,skip=4pt}
\setlength{\textfloatsep}{8pt plus 2pt minus 2pt}
\setlength{\floatsep}{6pt plus 2pt minus 2pt}
\setlength{\intextsep}{8pt plus 2pt minus 2pt}
\titleformat{\section}{\large\bfseries}{\thesection}{0.6em}{}
\titlespacing*{\section}{0pt}{1.1ex plus 0.3ex}{0.6ex}
\titleformat{\paragraph}[runin]{\bfseries}{}{0pt}{}[.~]
\titlespacing*{\paragraph}{0pt}{0.8ex}{0.5em}

\newcommand{\safeimg}[2][0.92\linewidth]{%
  \includegraphics[width=#1,height=0.38\textheight,keepaspectratio]{#2}%
}

\pagestyle{fancy}
\fancyhf{}
\fancyhead[L]{\footnotesize CISC3024 AI Assignment \#3}
\fancyhead[R]{\footnotesize CHE CHI HIN, Angus (UC325182)}
\fancyfoot[C]{\thepage}
\renewcommand{\headrulewidth}{0.3pt}

\title{\textbf{AI Assignment \#3}\\[0.25em]
\large Robust Sparse Smooth PCA (RSSPCA) for\\
Face Reconstruction and Recognition}
\author{CHE CHI HIN, Angus\\Student ID: UC325182\\CISC3024 Pattern Recognition}
\date{November 2026}

\begin{document}
\maketitle
\thispagestyle{fancy}

\begin{abstract}
\textbf{Thesis:} a 2025 PCA variant---Robust Sparse Smooth PCA (RSSPCA; Wang et al., PLOS One)---adds a graph-Laplacian smoothness constraint on top of L1 robustness and sparsity, and is a natural fit for classical face pattern recognition.
Under Cursor AI direction (no hand-written student code), we reimplemented RSSPCA in Python, compared it with PCA, PCA-L1, and RSPCA on the AT\&T ORL face database, and verified competitive clean reconstruction plus strong mid-dimensional recognition.
This report covers the required six sections: how AI found the algorithm, algorithm description, implementation, experiments, lessons learnt, and the source-code link.
\end{abstract}

\section{How I Asked AI Tools to Find the Algorithm}
The Moodle brief asks for a \emph{recent} PCA/LDA algorithm with application to computer vision and pattern recognition, with search, programming, and report writing done by an AI tool.
I used Cursor only.
I pointed at Assignments~1--2 and gave the agent this instruction:
\begin{quote}
\small
Find a recent PCA/LDA algorithm for computer vision / pattern recognition.
Do the search, the implementation, the experiments, and the report.
I will not write code.
Prefer a 2023--2025 method with a clear paper story (not another deep CNN like Assignment~1).
Ask clarifying questions, then lock one algorithm and a new independent repository like Assignment~1.
\end{quote}
The agent shortlisted four candidates (Table~\ref{tab:choice}).
I approved \textbf{RSSPCA} (Wang et al., PLOS One, May 2025)~\cite{wang2025rsspca}: it is clearly PCA-family, recent, has a public MATLAB reference, and targets face reconstruction and recognition.
Delivery matched Assignment~2: English LaTeX PDF, \texttt{process\_log.md}, and a public GitHub link.
Diary: \texttt{process\_log.md}.

\begin{table}[H]
\centering
\caption{Candidates considered for a recent PCA/LDA method that a laptop run can verify.}
\label{tab:choice}
\small
\begin{tabular}{lcccc}
\toprule
 & FF-PCA-LDA & 2DPCA+LDA & Robust Sparse LDA & RSSPCA \\
\midrule
Venue / year & Appl.\ Sci.\ 2022 & IAEAC 2024 & Pattern Recogn.\ 2024 & PLOS One 2025 \\
Family & PCA+LDA fusion & 2D PCA/LDA & LDA & PCA \\
CV task & plant disease & faces & images & faces \\
Public code & limited & limited & limited & yes (MATLAB) \\
Used here & no & no & no & yes \\
\bottomrule
\end{tabular}
\end{table}

\section{Algorithm Description}
\label{sec:algo}
\paragraph{From PCA to RSPCA}
Let training faces be mean-centered columns $X\in\mathbb{R}^{d\times n}$.
Classical PCA maximises $\|X^\top w\|_2^2$ under $\|w\|_2=1$.
PCA-L1 replaces the objective by $\|X^\top w\|_1$ for robustness to outliers~\cite{kwak2008pca}.
RSPCA further adds an L1 constraint on $w$ so loadings are sparse~\cite{meng2012rspca}.

\paragraph{RSSPCA objective}
RSSPCA keeps the L1 objective and adds a smoothness constraint via a pixel-grid Laplacian $L$~\cite{wang2025rsspca}:
\begin{equation}
\max_w \; \|X^\top w\|_1
\quad\text{s.t.}\quad
\|w\|_2^2=1,\;
\|w\|_1 \le c_1,\;
w^\top L w \le c_2.
\label{eq:rsspca}
\end{equation}
Here $L=D-A$ uses 4-neighbour adjacency on the $h\times w$ image lattice.
Then
\begin{equation}
w^\top L w = \tfrac12\sum_{i,j} A_{ij}(w_i-w_j)^2
\end{equation}
penalises dissimilar weights on adjacent pixels, encouraging spatially coherent ``eigenfaces''.

\paragraph{MM iterative solver}
Because of the L1 terms, the paper uses a Minorization--Maximization (MM) surrogate.
With sparsity / smoothness weights $\eta_1,\eta_2\ge 0$, each iteration (Algorithm~1 in the paper) does:
\begin{align}
U^{(k)}&=\mathrm{diag}(|w^{(k)}|),
\qquad
v^{(k)}=X\,\mathrm{sign}(X^\top w^{(k)}), \\
u^{(k)}&=(U^{(k)}+\eta_1 I+\eta_2 U^{(k)}L)^{-1} U^{(k)} v^{(k)}, \\
w^{(k+1)}&=u^{(k)}/\|u^{(k)}\|_2.
\end{align}
Multiple components use deflation $X\leftarrow (I-WW^\top)X$.
Our Python port adds Gram--Schmidt re-orthogonalisation so $W$ stays numerically orthonormal.
Intuition: $\eta_1$ trades sparsity vs.\ L2; $\eta_2$ trades smoothness vs.\ L2.
Setting $\eta_2=0$ recovers an RSPCA-like sparse robust PCA; $\eta_1=\eta_2=0$ recovers PCA-L1 behaviour.

\section{How AI Implemented the Algorithm}
The agent created a new project (not mixed into Assignment~1 FasterNet or Assignment~2 VBLL code):
\begin{verbatim}
CISC3024AI3/
  rsspca/           # Laplacian, PCA, PCA-L1, RSPCA, RSSPCA
  data/             # ORL loader (+ ensure_orl.py)
  experiments/      # reconstruction + recognition
  run_all.py
  outputs/          # figures + JSON metrics
  process_log.md
  report.tex
\end{verbatim}
Implementation choices:
\begin{itemize}
  \item \textbf{Language:} Python~3 + NumPy / SciPy sparse \texttt{spsolve} (readable vs.\ the MATLAB reference).
  \item \textbf{Baselines:} PCA (SVD), PCA-L1, RSPCA (soft-thresholding), RSSPCA (MM + Laplacian).
  \item \textbf{Data:} Cambridge AT\&T ORL (400 images, 40 subjects).
        sklearn Olivetti CDN returned HTTP 403; ORL is the same historical lab lineage.
        Images resized to $64\times 64$ ($d=4096$) for tractable solves.
  \item \textbf{Split:} 7 train / 3 test images per subject (seed 0).
  \item \textbf{Hyperparameters:} $\eta_1=0.001$, $\eta_2=0.1$, RSPCA $\lambda=0.05$, up to 40 components.
  \item \textbf{Tasks:} clean / block-occlusion / salt-and-pepper reconstruction (PSNR); 1-NN recognition in the projected subspace.
\end{itemize}
I did not write code myself; I approved the plan, checked figures/metrics, and submit the PDF on Moodle.

\section{Experiment Settings and Results}
\paragraph{Hardware / runtime}
CPU laptop; full \texttt{run\_all.py} $\approx$ 22--26\,s after data is cached.

\paragraph{Clean reconstruction}
Table~\ref{tab:psnr} reports PSNR (dB) of reconstructed test faces vs.\ \#components.
RSSPCA tracks classical PCA closely on clean faces (within $\sim$0.1\,dB at $r=40$), confirming the MM solver and deflation are stable after re-orthogonalisation.

\begin{table}[H]
\centering
\caption{Clean reconstruction PSNR (dB) on ORL test faces.}
\label{tab:psnr}
\small
\begin{tabular}{lcccc}
\toprule
\#comp & PCA & PCA-L1 & RSPCA & RSSPCA \\
\midrule
5  & 19.61 & 19.55 & 19.55 & 19.55 \\
10 & 20.89 & 20.78 & 20.78 & 20.78 \\
20 & 22.04 & 21.91 & 21.91 & 21.94 \\
30 & 22.75 & 22.61 & 22.60 & 22.65 \\
40 & 23.31 & 23.16 & 23.15 & 23.21 \\
\bottomrule
\end{tabular}
\end{table}

\paragraph{Recognition (1-NN)}
Table~\ref{tab:recog} and Figure~\ref{fig:recog} report nearest-neighbour accuracy in the projected subspace.
At compact dimensions ($r=5,10$), RSSPCA is best or tied-best---useful when one wants few features.
At very high $r$, plain PCA slightly leads on this clean ORL split (ceiling effect; all methods $>96\%$).

\begin{table}[H]
\centering
\caption{1-NN recognition accuracy on ORL (7 train / 3 test per subject).}
\label{tab:recog}
\small
\begin{tabular}{lcccc}
\toprule
\#comp & PCA & PCA-L1 & RSPCA & RSSPCA \\
\midrule
5  & 0.842 & 0.842 & 0.842 & \textbf{0.850} \\
10 & 0.950 & 0.958 & 0.958 & \textbf{0.967} \\
15 & 0.950 & 0.942 & 0.942 & 0.950 \\
20 & 0.950 & 0.950 & 0.950 & 0.950 \\
30 & 0.958 & 0.958 & 0.958 & 0.958 \\
40 & \textbf{0.975} & 0.967 & 0.967 & 0.967 \\
\bottomrule
\end{tabular}
\end{table}

\begin{figure}[H]
\centering
\safeimg[0.88\linewidth]{recognition_accuracy_curves.png}
\caption{1-NN accuracy vs.\ number of components on ORL.}
\label{fig:recog}
\end{figure}

\begin{figure}[H]
\centering
\safeimg[0.88\linewidth]{reconstruction_psnr_curves.png}
\caption{Clean and occluded reconstruction PSNR vs.\ \#components.}
\label{fig:psnr}
\end{figure}

\paragraph{Robustness probes}
On $16\times 16$ block occlusion (project occluded test, score vs.\ clean GT, $r=40$): PCA 18.88\,dB, PCA-L1 18.82, RSPCA 18.81, RSSPCA 18.80---methods stay close; occlusion geometry is structured, so gains are modest.
On 10\% salt-and-pepper ($r=30$): PCA 22.32, PCA-L1 22.20, RSPCA 22.20, RSSPCA 22.24\,dB (Figure~\ref{fig:sp}).

\begin{figure}[H]
\centering
\safeimg[0.72\linewidth]{saltpepper_psnr_bars.png}
\caption{Salt-and-pepper (10\%) reconstruction PSNR at $r=30$.}
\label{fig:sp}
\end{figure}

\paragraph{Qualitative eigenfaces}
Figure~\ref{fig:eigen} compares learned components.
RSSPCA loadings are spatially smoother than raw sparse RSPCA while remaining more localised than dense PCA.
Reconstruction galleries are in Figures~\ref{fig:galc}--\ref{fig:galo}.

\begin{figure}[H]
\centering
\begin{minipage}{0.48\linewidth}
\centering
\safeimg[0.98\linewidth]{eigenfaces_pca.png}
\end{minipage}\hfill
\begin{minipage}{0.48\linewidth}
\centering
\safeimg[0.98\linewidth]{eigenfaces_rsspca.png}
\end{minipage}
\caption{Left: PCA components. Right: RSSPCA components (smooth + sparse structure).}
\label{fig:eigen}
\end{figure}

\begin{figure}[H]
\centering
\safeimg[0.95\linewidth]{recon_gallery_clean.png}
\caption{Clean reconstruction gallery ($r=30$).}
\label{fig:galc}
\end{figure}

\begin{figure}[H]
\centering
\safeimg[0.95\linewidth]{recon_gallery_occluded.png}
\caption{Block-occlusion input (top) and reconstructions ($r=30$).}
\label{fig:galo}
\end{figure}

\section{What I Have Learnt from This AI Assignment}
\begin{itemize}
  \item Prompting beats typing: locking ``recent PCA/LDA + CV + AI does all code'' early steered the agent away from another CNN Assignment~1 clone.
  \item Recent $\neq$ automatically better on every metric: on clean ORL, classical PCA remains a strong baseline at high dimension; RSSPCA's value shows up as a \emph{principled} robust/sparse/smooth model and better mid-$r$ accuracy.
  \item Numerical details matter: without Gram--Schmidt, deflation $(I-WW^\top)X$ became unstable and RSSPCA PSNR collapsed after $\sim$20 components---an implementation lesson, not a paper flaw.
  \item Data logistics are part of the experiment: figshare Olivetti 403 forced an ORL fallback; documenting that is honest science.
  \item Process logs help Turnitin-facing reports: the six required sections map cleanly when search $\rightarrow$ math $\rightarrow$ code $\rightarrow$ results are recorded chronologically.
\end{itemize}

\section{Source Code Webpage Link}
Source code:\\[0.3em]
\url{https://github.com/AngusJai/CISC3024-AI-Assignment3}\\[0.3em]
The repository contains the solvers, experiment scripts, \texttt{process\_log.md}, \texttt{report.tex}, and the metrics and figures under \texttt{outputs/}.
Dataset downloads and virtual environments are gitignored; \texttt{att\_faces.zip} is included for reproducibility.
Reproduce with:
\begin{verbatim}
python -m venv .venv && source .venv/bin/activate
pip install -r requirements.txt
python scripts/ensure_orl.py
python run_all.py
\end{verbatim}

\section{Conclusion}
RSSPCA is a recent PCA algorithm that an agent can reimplement from a published MM update and test on face images in one afternoon.
With $(\eta_1,\eta_2)=(0.001,0.1)$ it matches PCA on clean reconstruction and leads at compact recognition dimensions ($r=5,10$).
At high $r$ on clean ORL, classical PCA remains slightly ahead---a ceiling effect, not a reason to discard the method.
The assignment value is the end-to-end AI workflow: literature search, correct sparse linear algebra, honest tables, and a submit-ready report without student-written code.

\begin{thebibliography}{9}
\bibitem{wang2025rsspca}
J.~Wang, X.~Xie, L.~Zhang, J.~Li, H.~Cai, and Y.~Feng.
Robust sparse smooth principal component analysis for face reconstruction and recognition.
PLOS One, 20(5):e0323281, 2025.
\bibitem{kwak2008pca}
N.~Kwak.
Principal component analysis based on L1-norm maximization.
IEEE Transactions on Pattern Analysis and Machine Intelligence, 2008.
\bibitem{meng2012rspca}
D.~Meng, Q.~Zhao, and Z.~Xu.
Improve robustness of sparse PCA by L1-norm maximization.
Pattern Recognition, 2012.
\bibitem{samaria1994orl}
F.~S.~Samaria and A.~C.~Harter.
Parameterisation of a stochastic model for face identification.
IEEE Workshop on Applications of Computer Vision, 1994.
\end{thebibliography}

\end{document}
